\section{Kiến trúc hệ thống}

\subsection{Bản đồ codebase}

\begin{lstlisting}[style=rae]
configs/
  stage1/
  stage2/
src/
  stage1/
  stage2/
  utils/
  train.py
  sample.py
  sample_ddp.py
  stage1_sample.py
  stage1_sample_ddp.py
  build_fid_stats.py
  evaluate_fid.py
src_jax/
  vendor.py
  config_adapter.py
  stage2_runtime.py
  stage1_runtime.py
  build_stage1_stats.py
  export_celebahq_hf.py
  export_celebahq_tfds.py
  reconstruct_folder.py
  hf_utils.py
  train.py
  sample.py
  sample_ddp.py
  stage1_sample.py
  push_hf.py
vaes-jax-celebahq-kaggle.ipynb
vaes-jax-celebahq-kaggle-tpuv5e8-sitb.ipynb
vaes-jax-celebahq-kaggle-tpuv5e8-sitb-resume.ipynb
\end{lstlisting}

\subsection{Stage 1: Representation Autoencoder và StabilityVAE}

Thành phần gốc của Stage 1 nằm trong \texttt{src/stage1/rae.py}. Lớp
\texttt{RAE} kết hợp:

\begin{itemize}[leftmargin=1.5em]
  \item encoder ở \texttt{src/stage1/encoders/},
  \item decoder ở \texttt{src/stage1/decoders/},
  \item latent normalization qua \texttt{normalization\_stat\_path},
  \item latent noising qua \texttt{noise\_tau}.
\end{itemize}

Hành vi chính:

\begin{itemize}[leftmargin=1.5em]
  \item \texttt{encode(x)} resize đầu vào về đúng độ phân giải của encoder,
  chuẩn hóa theo image processor, chạy encoder, reshape latent về dạng 2D nếu
  cần, rồi áp normalization;
 \item \texttt{decode(z)} đảo normalization, đổi latent về dạng sequence nếu
  cần, rồi giải mã thành ảnh RGB bằng decoder ViT.
\end{itemize}

Ở nhánh này, \texttt{src/stage1/stability\_vae.py} còn cung cấp target
\texttt{stage1.StabilityVAE} như một alias repo-facing cho encoder
\texttt{StabilityVAE} native của backend NNX. Đường này:

\begin{itemize}[leftmargin=1.5em]
  \item không cần \texttt{pretrained\_decoder\_path} của RAE;
  \item không bắt buộc phải bootstrap latent stat trước khi chạy CelebA-HQ;
  \item mặc định suy ra latent shape \texttt{[4, 32, 32]} ở độ phân giải
  \texttt{256x256}.
\end{itemize}

\subsection{Stage 2: SiT trong không gian latent}

Surface Stage 2 dùng chung của branch này nằm ở \texttt{src/stage2/models/SiT.py}.
File này giờ có hai alias repo-facing:

\begin{itemize}[leftmargin=1.5em]
  \item \texttt{SiTDH}: backbone DH/two-tower dựa trên
  \texttt{DiTwDDTHead};
  \item \texttt{SiT}: backbone single-tower dựa trên
  \texttt{LightningDiT}.
\end{itemize}

Nhờ đó nhánh hiện tại vẫn giữ được đường \texttt{SiTDH} ở mức code/config khi
cần, đồng thời public notebook flow của CelebA-HQ tập trung vào
\texttt{StabilityVAE + SiT-B}. Objective trên đường JAX vẫn đi qua interface
\texttt{sit} cho cả hai dạng backbone.

Các đặc điểm quan trọng:

\begin{itemize}[leftmargin=1.5em]
  \item đầu vào là latent thay vì pixel;
  \item embedding thời gian bằng Gaussian Fourier features;
  \item class conditioning bằng \texttt{LabelEmbedder};
  \item tùy chọn ROPE, RMSNorm, SwiGLU, qk-norm;
  \item hỗ trợ classifier-free guidance và autoguidance khi sampling.
\end{itemize}

Đường JAX ánh xạ \texttt{stage\_2.target} sang backbone tương ứng của backend
NNX:

\begin{itemize}[leftmargin=1.5em]
  \item \texttt{SiTDH} $\rightarrow$ \texttt{lightning\_ddt},
  \item \texttt{SiT} $\rightarrow$ \texttt{lightning\_dit},
  \item \texttt{LightningDiT} $\rightarrow$ \texttt{lightning\_dit},
  \item các target DDT legacy $\rightarrow$ \texttt{lightning\_ddt},
  \item các target kiểu DiT khác $\rightarrow$ \texttt{dit}.
\end{itemize}

\subsection{Transport objective}

Phần này nằm trong \texttt{src/stage2/transport/transport.py}. Đây là lớp trung
gian giữa latent diffusion model và train loop.

Nhiệm vụ của transport:

\begin{itemize}[leftmargin=1.5em]
  \item lấy mẫu thời gian \texttt{t};
  \item xây dựng đường nối giữa noise và dữ liệu thật;
  \item tính target phù hợp với loại output của model;
  \item trả về drift function để phục vụ ODE hoặc SDE sampler.
\end{itemize}

Ở đường JAX, adapter không giữ lại implementation transport gốc này. Nó ánh xạ
các ý nghĩa tương đương sang interface của backend NNX, chủ yếu là SiT với
sampling kiểu Euler.

\subsection{Lớp adapter JAX}

\begin{center}
\begin{longtable}{>{\raggedright\arraybackslash}p{0.28\linewidth} >{\raggedright\arraybackslash}p{0.62\linewidth}}
\toprule
File & Vai trò \\
\midrule
\endhead
\texttt{src\_jax/vendor.py} & Bootstrap \texttt{diffuse\_nnx} vào \texttt{\~{}/.cache/rae\_jax/diffuse\_nnx} và pin commit \\
\texttt{src\_jax/config\_adapter.py} & Chuyển YAML OmegaConf của repo sang config backend NNX, map \texttt{SiT}/\texttt{SiTDH}/\texttt{StabilityVAE} sang backend phù hợp, và forward \texttt{random\_flip} / prefetch knobs \\
\texttt{src\_jax/stage2\_runtime.py} & Train, sampling, load checkpoint PyTorch hoặc Orbax, JAX validation-loss integration, FID glue, wandb bridge \\
 \texttt{src\_jax/stage1\_runtime.py} & Load encoder Stage 1 JAX, reconstruction một ảnh, reconstruction cả thư mục và tích lũy latent stat cho cả \texttt{RAE} lẫn \texttt{StabilityVAE} \\
 \texttt{src\_jax/build\_stage1\_stats.py} & Tạo file \texttt{stat.pt} theo định dạng mean / var để dùng lại cho Stage 1 \\
 \texttt{src\_jax/export\_celebahq\_hf.py} & Download \texttt{eurecom-ds/celeba-hq-256} từ Hugging Face rồi export thành cây \texttt{ImageFolder} cho pipeline JAX hiện tại \\
 \texttt{src\_jax/export\_celebahq\_tfds.py} & Export TFDS \texttt{celeb\_a\_hq/256} thành cây \texttt{ImageFolder} cho pipeline JAX hiện tại và ép protobuf runtime kiểu Python để tránh lỗi import TFDS trên Kaggle \\
 \texttt{src\_jax/build\_fid\_stats.py} & Build \texttt{fid\_ref} bằng đúng detector Flax Inception của backend JAX \\
 \texttt{src\_jax/reconstruct\_folder.py} & Export reconstruction cho cả thư mục ảnh theo pipeline JAX \\
\texttt{src\_jax/hf\_utils.py} & Upload file hoặc thư mục lên Hugging Face \\
\bottomrule
\end{longtable}
\end{center}

\subsection{Hai mô hình thực thi}

\subsubsection{TorchXLA}

Các entrypoint distributed trong \texttt{src/} dùng \texttt{xmp.spawn(...)} và
mỗi TPU core chạy một process riêng. Các primitive quan trọng:

\begin{itemize}[leftmargin=1.5em]
  \item \texttt{DistributedSampler},
  \item \texttt{ParallelLoader},
  \item \texttt{xm.optimizer\_step(...)},
  \item \texttt{xm.rendezvous(...)} và \texttt{xm.mesh\_reduce(...)}.
\end{itemize}

\subsubsection{JAX / NNX}

Đường \texttt{src\_jax/} không tự quản lý toàn bộ graph, checkpoint và sampler ở
mức thấp. Nó dựa vào backend NNX để:

\begin{itemize}[leftmargin=1.5em]
  \item tạo model, optimizer, sampler và EMA,
  \item quản lý checkpoint Orbax,
  \item chạy sampling và FID trong định dạng tensor NHWC của JAX,
  \item mở rộng sang nhiều process JAX nếu môi trường đã cấu hình phù hợp.
\end{itemize}

Repo cũng vá backend này để EMA khởi tạo từ bản copy của model hiện thời thay vì
một cây tham số toàn số 0. Nhờ vậy \texttt{ema\_network\_samples} và online FID
không còn bắt đầu từ một EMA cold-start sai ngay từ step đầu.

\subsection{Checkpoint và tương thích trọng số}

Với branch \texttt{jax-vae-sit-celebahq256}, checkpoint Stage 2 phải khớp shape
của backbone đã chọn:

\begin{itemize}[leftmargin=1.5em]
  \item nếu \texttt{stage\_2.ckpt} là file PyTorch \texttt{.pt}, adapter sẽ
  port state dict sang NNX để inference hoặc khởi tạo train nếu checkpoint đó
  đã tương thích với \texttt{SiT} hoặc \texttt{SiTDH};
  \item nếu \texttt{stage\_2.ckpt} là thư mục checkpoint Orbax, adapter sẽ
  restore trực tiếp run JAX trước đó với cùng shape backend;
  \item checkpoint DDT legacy không được tự động convert trên branch này;
  \item decoder Stage 1 ở đường JAX vẫn dùng được checkpoint decoder PyTorch
  trên flow RAE; flow \texttt{StabilityVAE} dùng trọng số VAE native của backend.
\end{itemize}
